\section*{Capítulo \ref{ch:b2:acyl-substitution} --- Derivados dos ácidos carboxílicos}

\begin{solution}{exo:b2:acyl-substitution:1}
Etanamida $<$ etanoato de etila $<$ anidrido etanoico $<$ cloreto de etanoíla.
\end{solution}

\begin{solution}{exo:b2:acyl-substitution:2}
Ácido etanoico (+ \ce{HCl}); etanoato de etila; etanamida (+ cloreto de amônio, o segundo
equivalente de amônia capturando o \ce{HCl}); anidrido etanoico (+ \ce{NaCl});
\textit{N},\textit{N}-dimetiletanamida (+ cloreto de trietilamônio).
\end{solution}

\begin{solution}{exo:b2:acyl-substitution:3}
\ce{C6H5COOCH3 + NaOH -> C6H5COONa + CH3OH}. $13.6/136.15 = \qty{0.0999}{mol}$, a mesma quantidade
de \ce{NaOH}: \qty{4.00}{g}.
\end{solution}

\begin{solution}{exo:b2:acyl-substitution:4}
$\gamma$-butirolactona (oxolan-2-ona): um anel de cinco membros, quatro carbonos e um oxigênio.
\end{solution}

\begin{solution}{exo:b2:acyl-substitution:5}
A amônia se adiciona ao carbono carbonílico; um próton passa do nitrogênio para o oxigênio (ou
para o grupo de saída); o etóxido (como etanol, após transferência de próton) sai e a ligação
\ce{C=O} se refaz: etanamida e etanol.
\end{solution}

\begin{solution}{exo:b2:acyl-substitution:6}
O \ce{OH} fenólico é acilado (catálise ácida): ácido 2-acetoxibenzoico; subproduto, o ácido etanoico.
\end{solution}

\begin{solution}{exo:b2:acyl-substitution:7}
2-metilpropan-2-ol. A primeira adição dá a propanona, mais eletrofílica que o éster, que
recebe de imediato a segunda metila.
\end{solution}

\begin{solution}{exo:b2:acyl-substitution:8}
Butano-1,4-diol: o éster é reduzido a dois álcoois primários, com abertura do anel.
\end{solution}

\begin{solution}{exo:b2:acyl-substitution:9}
Um equivalente: $x^2/(1 - x)^2 = 4$, $x = 2/3$. Três equivalentes: $x^2/[(1 - x)(3 - x)] = 4$, $3x^2 - 16x
+ 12 = 0$, $x = 0.90$. Remover a água à medida que se forma (Dean--Stark) a desloca ainda mais.
\end{solution}

\begin{solution}{exo:b2:acyl-substitution:10}
O \ce{SOCl2} dá o cloreto de benzoíla; com dois equivalentes de dietilamina (ou um mais
trietilamina) ele dá a amida. O ácido e a amina misturados diretamente sofrem uma reação
ácido--base: benzoato de dietilamônio, que só dá a amida sob aquecimento forte.
\end{solution}

\begin{solution}{exo:b2:acyl-substitution:11}
Três KOH por trioleína: $3 \times 56.11/885.4 = \qty{0.190}{g}$ por grama, um índice de \omterm{def:b2:acyl-substitution:saponification}{saponificação}
de 190.
\end{solution}

\begin{solution}{exo:b2:acyl-substitution:12}
\ce{NaBH4}: só a cetona, a álcool secundário. \ce{LiAlH4}: a cetona ao álcool secundário, o éster
a um álcool primário (liberando a parte álcool), a amida a uma amina.
\end{solution}

\begin{solution}{pb:b2:acyl-substitution:1}
\textbf{1.} Ácido oleico \ce{C18H34O2}; glicerol \ce{C3H8O3}.
\textbf{2.} \ce{C57H104O6}: $57 \times 12.011 + 104 \times 1.008 + 6 \times 15.999 = \qty{885.45}{g/mol}$.
\textbf{3.} Três.
\textbf{4.} A ligação dupla cis dobra as cadeias, que se empacotam mal: forças intermoleculares
fracas, um ponto de fusão baixo.
\textbf{5.} Glicerol e três oleatos de sódio, um sabão.
\textbf{6.} Uma mistura de ésteres metílicos de ácidos graxos.
\textbf{7.} \ce{C57H104O6 + 3CH3OH -> C3H8O3 + 3C19H36O2}.
\textbf{8.} O íon metóxido, formado a partir do metanol e da base, é um nucleófilo muito mais forte
que o metanol; ele ataca as carbonilas dos ésteres.
\textbf{9.} O carbono carbonílico de um grupo éster, levando \ce{O-}, o \ce{OCH3} do metóxido e o
oxigênio do grupo glicerila: tetraédrico.
\textbf{10.} A \omterm{def:b2:acyl-substitution:transesterification}{transesterificação} é um equilíbrio com constante perto de 1; o excesso de metanol o
desloca para os ésteres metílicos.
\textbf{11.} O glicerol se separa em uma camada inferior: remover um produto desloca o equilíbrio.
\textbf{12.} Não: o metóxido é regenerado a cada troca.
\textbf{13.} Ele neutraliza a base, dando oleato de sódio (sabão): o catalisador é consumido.
\textbf{14.} Ela os saponifica: o hidróxido (vindo da água e do metóxido) corta os ésteres em sabão e
metanol, de modo irreversível.
\textbf{15.} Ele consome base, emulsiona a mistura e impede a separação da camada de glicerol.
\textbf{16.} Por uma \omterm{def:b2:acyl-substitution:esterification}{esterificação} prévia com metanol, catalisada por ácido, que transforma os ácidos
livres em ésteres metílicos.
\textbf{17.} A água hidrolisa os ésteres e converte o catalisador em hidróxido, que saponifica.
\textbf{18.} $10^6/885.45 = \qty{1129}{mol}$.
\textbf{19.} $3 \times 1129 \times 32.04 = \qty{108.6}{kg}$.
\textbf{20.} $3 \times 1129 \times 296.50 = \qty{1004.6}{kg}$.
\textbf{21.} $1000 + 108.6 = \qty{1108.6}{kg}$ na entrada; $104.0 + 1004.6 = \qty{1108.6}{kg}$ na saída.
\textbf{22.} \qty{20}{kg} de ácido oleico, $20\,000/282.47 = \qty{70.8}{mol}$, a mesma quantidade de \ce{NaOH}:
\qty{2.83}{kg}.
\textbf{23.} $1129 \times 92.09 = \boldsymbol{\approx \qty{104}{kg}}$ de glicerol por tonelada.
\end{solution}
